\section{Income and wealth concepts}

%% CHANGE: Added proper \citep for Miller-Modigliani; noted benchmark nature of M&M assumption
National income, as defined by the Bureau of Economic Analysis, consists of the total labor and capital income earned in production. This definition explicitly excludes capital gains (which are not earned by production). Indirectly, however, capital gains are present through the retained earnings of corporations. Under the \citet{miller1961dividend} irrelevance assumption, a firm that retains a dollar of earnings increases its market value by \$1. In this benchmark case, at least some capital gain income is accounted for in the national accounts income measure.

%% CHANGE: Fixed typos "fows"->"flows", "defne"->"define"; clarified "net of contributions"; improved "pure capital gains" definition
To isolate capital gains from other income flows, we define ordinary factor income as national income minus the retained earnings of corporations. Capital gains for year $t$ are the real increase in the market value of an individual's financial and nonfinancial assets over the period, net of new contributions and purchases. Pure capital gains equal total capital gains minus retained earnings---the component of appreciation that is not accounted for in national income.

%% CHANGE: Specified which capital income sources; added total (pure + RE) figure
Figure 2 compares capital gains to dividends, interest, and rental income in the post-war era. Pure capital gains are large, especially in the post-1980 period, where they average 10\% of national income. Retained earnings, the component present in national income, averaged 4\% of national income for the same period. Total capital gains---including both pure gains and retained earnings---thus averaged 14\% of national income since 1980.

%% CHANGE: Added formal citations for Haig, Simons, Hicks; added displayed equation; stated equivalence with NI + pure KG
Haig-Simons income equals ordinary factor income plus capital gains. The concept originates with \citet{haig1921concept} and \citet{simons1938personal}. As described by \citet{hicks1939value}, Haig-Simons income ``is what [one] can consume during the week and still expect to be as well off at the end of the week as at the beginning.'' Formally, Haig-Simons income for individual $i$ in year $t$ is:
\be
Y^{HS}_{it} = Y^{factor}_{it} + KG_{it} = C_{it} + \Delta W_{it}
\ee
where $Y^{factor}_{it}$ is ordinary factor income, $KG_{it}$ is total capital gains, $C_{it}$ is consumption, and $\Delta W_{it}$ is the change in net worth. Equivalently, since ordinary factor income equals national income minus retained earnings, Haig-Simons income equals national income plus pure capital gains.

%% CHANGE: Removed dangling "Existinig empirical work" fragment; fixed "ceteris paribus"; restructured to lead with practical case; added transition to discount rate concern
We include capital gains as part of an income measure because they expand the stock of resources individuals have available to consume, save, or pass on as bequests. Capital gains enable higher spending directly through asset sales and indirectly by relaxing collateral and borrowing constraints. From a theoretical perspective, ceteris paribus (importantly, holding constant expected returns), an increase in the value of an individual's assets increases their utility by the value of the capital gains times the marginal utility of wealth \citep{cochrane2011discount}. However, this reasoning assumes that capital gains reflect changes in fundamentals; gains driven by discount rate changes may not increase welfare proportionally.

%% CHANGE: Fixed awkward "More generally than"; integrated Fagereng/Aguiar result from comments into main text; added Auclert citation
An important case of capital gains not driven by fundamentals is the secular decline in real interest rates since the 1980s, both in the United States \citep{eggertsson2019model} and globally \citep{rachel2019secular}. Lower interest rates generate higher asset prices (and positive capital gains), since future dividends are discounted at lower rates. \citet{auclert2019monetary} shows that for a temporary change in interest rates, the welfare effect is proportional to a household's unhedged interest rate exposure, which may differ substantially from the capital gain. More generally, asset prices may be driven by variations in \textit{discount rates}, defined as anything that changes the price of an asset without affecting its present or future cash flow. \citet{fagereng2025asset} and \citet{aguiar2024putting} show that discount rates affect welfare both through capital gains and through changing effective future returns. Under a number of assumptions,\footnote{Such as: asset price changes are infinitesimally small and thus do not affect asset allocations, asset price changes do not relax borrowing constraints, there is no change in the underlying risk of the asset, and wealth does not appear in the utility function.} the change in welfare from a discount rate shift is proportional to the present value of an individual's net asset sales, rather than their capital gains.

%% CHANGE: Fixed typos ("earniings", "discount changes in the discount rate"); fixed "five-year moving average"; changed "decline" to "changes" for symmetry; added framing sentence; added brief explanation of why smoothing helps
We address this concern through three complementary approaches in Section NNN. First, we estimate asset pricing models for public equities, private business, and housing wealth over our sample period, and strip out from the capital gain series any asset price changes driven by changes in interest rates. Second, for our main capital gain results we smooth out short- and medium-term discount rate variation by using a five-year moving average of returns, and also show results for a 10-year smoothed average; because short-run discount rate variation tends to mean-revert, multi-year averages better capture the permanent component of capital gains. Finally, we present evidence that, in line with existing literature, the majority of long-run variation in asset prices is due to earnings changes and not changes in the discount rate.

%% CHANGE: Clarified pension/DI/UI exception; "networth"->"net worth"; noted C+DeltaW is the definition not approximation; added transition for alternative measures list
At the individual level, income differs based on whether it is pre- or post-tax. ``Pre-tax'' Haig-Simons income, following PSZ, is income before taxes and transfers, but inclusive of pension distributions, disability insurance, and unemployment insurance. ``Post-tax'' Haig-Simons income equals pre-tax income, minus taxes, plus all forms of government spending: cash transfers, in-kind transfers, and collective consumption expenditures. Pre-tax Haig-Simons income is useful for measuring the effective tax rate faced by individuals, while post-tax Haig-Simons income most closely corresponds to the textbook definition $C_t + \Delta W_t$. To facilitate comparison with prior work, we also explore the robustness of our results by adding our capital gains series to other income measures: the \textit{market income} of PSZ, the comprehensive income of \citet{larrimore2021recent}, expanded income of the IRS, and AGI minus realized capital gains.
